% Manual de Usuario — NetManagement
% Documento autocontenido para compilar de forma independiente
\documentclass[12pt,a4paper]{article}

% ── Paquetes ──────────────────────────────────────────────────
\usepackage[utf8]{inputenc}
\usepackage[T1]{fontenc}
\usepackage[spanish]{babel}
\usepackage{geometry}
\geometry{margin=2.5cm}
\setlength{\headheight}{14pt}
\usepackage{graphicx}
\usepackage{hyperref}
\hypersetup{
    colorlinks=true,
    linkcolor=blue!70!black,
    urlcolor=blue!60!black,
    citecolor=green!50!black,
}
\usepackage{xcolor}
\usepackage{tikz}
\usepackage{booktabs}
\usepackage{tabularx}
\usepackage{enumitem}
\usepackage{float}
\usepackage{fancyhdr}
\usepackage{titlesec}
\usepackage{mdframed}

% ── Colores ───────────────────────────────────────────────────
\definecolor{primary}{HTML}{1565C0}
\definecolor{secondary}{HTML}{0D47A1}
\definecolor{accent}{HTML}{FF6F00}
\definecolor{notebg}{HTML}{E3F2FD}
\definecolor{noteframe}{HTML}{1565C0}
\definecolor{warnbg}{HTML}{FFF8E1}
\definecolor{warnframe}{HTML}{FF6F00}
\definecolor{tipbg}{HTML}{E8F5E9}
\definecolor{tipframe}{HTML}{2E7D32}

% ── Cajas informativas ───────────────────────────────────────
\newmdenv[
    backgroundcolor=notebg,
    linecolor=noteframe,
    linewidth=3pt,
    topline=false, bottomline=false, rightline=false,
    innerleftmargin=12pt, innerrightmargin=8pt,
    innertopmargin=6pt, innerbottommargin=6pt,
]{infonote}

\newmdenv[
    backgroundcolor=warnbg,
    linecolor=warnframe,
    linewidth=3pt,
    topline=false, bottomline=false, rightline=false,
    innerleftmargin=12pt, innerrightmargin=8pt,
    innertopmargin=6pt, innerbottommargin=6pt,
]{warnnote}

\newmdenv[
    backgroundcolor=tipbg,
    linecolor=tipframe,
    linewidth=3pt,
    topline=false, bottomline=false, rightline=false,
    innerleftmargin=12pt, innerrightmargin=8pt,
    innertopmargin=6pt, innerbottommargin=6pt,
]{tipnote}

% ── Encabezados y pies ───────────────────────────────────────
\pagestyle{fancy}
\fancyhf{}
\fancyhead[L]{\small\textit{Manual de Usuario --- NetManagement}}
\fancyhead[R]{\small\thepage}
\renewcommand{\headrulewidth}{0.4pt}

% ── Secciones ─────────────────────────────────────────────────
\titleformat{\section}{\Large\bfseries\color{primary}}{}{0em}{\thesection\quad}
\titleformat{\subsection}{\large\bfseries\color{secondary}}{}{0em}{\thesubsection\quad}

% ==============================================================
\begin{document}

% ── Portada ───────────────────────────────────────────────────
\begin{titlepage}
  \thispagestyle{empty}
  \centering

  \vspace*{1cm}

  {\large\bfseries\color{secondary} IES San Mamede\par}
  \vspace{0.3cm}
  {\normalsize Ciclo Formativo de Grado Superior\\[0.15cm]
  \textbf{Desarrollo de Aplicaciones Multiplataforma (DAM)}\par}

  \vspace{1.5cm}

  {\Huge\bfseries\color{primary} NetManagement\par}
  \vspace{0.6cm}
  {\LARGE Sistema de Monitorización y Control\\de Equipos en Aulas Informáticas\par}

  \vspace{1.5cm}

  \begin{tikzpicture}
    \draw[primary, line width=2pt] (0,0) -- (12,0);
  \end{tikzpicture}

  \vspace{1.5cm}

  {\Large\bfseries Manual de Usuario\par}

  \vspace{2cm}

  \begin{tabular}{rl}
    \textbf{Autor:}  & Enooc Domínguez Quiroga \\[0.3cm]
    \textbf{Centro:} & IES San Mamede \\[0.3cm]
    \textbf{Curso:}  & 2025 / 2026 \\[0.3cm]
    \textbf{Fecha:}  & Junio de 2026 \\
  \end{tabular}

  \vfill

  {\small Versión 1.0 --- Junio 2026\par}
\end{titlepage}

\newpage
\thispagestyle{empty}
\tableofcontents
\newpage

% ==============================================================
\section{Introducción}
% ==============================================================

\textbf{NetManagement} es una aplicación web para la monitorización y control de equipos informáticos en aulas. Permite al profesorado y al personal de administración visualizar el estado de los equipos en tiempo real, controlar el acceso a internet de forma individual o por aula completa, y gestionar los recursos de red del centro.

Este manual describe el uso de la interfaz web desde el punto de vista del usuario. No cubre la instalación ni la configuración del servidor (véase el \textit{Manual de Instalación}).

\subsection{Tipos de usuario}

La aplicación distingue tres niveles de acceso:

\begin{table}[H]
\centering
\begin{tabularx}{\textwidth}{lX}
    \toprule
    \textbf{Tipo} & \textbf{Capacidades} \\
    \midrule
    \textbf{Usuario estándar} & Accede únicamente al Dashboard. Visualiza los dispositivos de las aulas que tiene asignadas. No puede acceder a secciones de administración. \\[0.3cm]
    \textbf{Staff (profesor / técnico)} & Acceso completo al Dashboard y a todas las secciones de gestión: aulas, equipos de red, hosts permitidos, usuarios y grupos. \\[0.3cm]
    \textbf{Superusuario} & Acceso total sin restricciones. Puede ver dispositivos de cualquier aula. Tiene todos los permisos de staff más la capacidad de gestionar otros superusuarios. \\
    \bottomrule
\end{tabularx}
\end{table}

\begin{infonote}
\textbf{Nota:} Los permisos también pueden refinarse mediante grupos de permisos. Un usuario de staff puede tener acceso solo a determinadas secciones según el grupo al que pertenezca.
\end{infonote}

% ==============================================================
\section{Acceso a la aplicación}
% ==============================================================

Abra su navegador web y acceda a la URL proporcionada por el administrador del sistema (p. ej., \texttt{https://netmanagement.ejemplo.es}).

\subsection{Inicio de sesión}

En la pantalla de inicio de sesión:

\begin{enumerate}
    \item Introduzca su \textbf{nombre de usuario} en el campo «Usuario».
    \item Introduzca su \textbf{contraseña} en el campo «Contraseña».
    \item Pulse el botón \textbf{Iniciar sesión}.
\end{enumerate}

Si las credenciales son correctas, será redirigido al \textbf{Dashboard} (página de inicio).

\begin{warnnote}
\textbf{Atención:} Si aparece el mensaje «Credenciales inválidas», compruebe que el usuario existe y que la contraseña es correcta. Contacte con el administrador si no recuerda sus credenciales.
\end{warnnote}

\subsection{Cierre de sesión}

Para cerrar sesión de forma segura, haga clic en \textbf{Cerrar sesión} en la barra lateral izquierda (icono de salida, en la parte inferior). La sesión se invalida y se le redirige a la página de inicio de sesión.

% ==============================================================
\section{Navegación general}
% ==============================================================

Tras iniciar sesión, la interfaz se divide en dos zonas principales:

\begin{itemize}
    \item \textbf{Barra lateral izquierda:} Menú de navegación con acceso a todas las secciones disponibles según su nivel de acceso.
    \item \textbf{Área de contenido:} Muestra la sección actualmente seleccionada.
\end{itemize}

\subsection{Opciones del menú lateral}

\begin{table}[H]
\centering
\begin{tabularx}{\textwidth}{llX}
    \toprule
    \textbf{Icono} & \textbf{Sección} & \textbf{Descripción} \\
    \midrule
    Dashboard      & Dashboard            & Vista principal con los dispositivos del aula seleccionada. \\
    Classroom      & Aulas                & Gestión de aulas y asignación de dispositivos. \textit{(solo staff)} \\
    Switches       & Switches / Routers   & Gestión de equipos de red MikroTik. \textit{(solo staff)} \\
    Hosts          & Hosts Permitidos     & Dominios o IPs con acceso permitido durante el bloqueo. \textit{(solo staff)} \\
    Personas       & Usuarios             & Gestión de cuentas de usuario. \textit{(solo staff)} \\
    Grupos         & Grupos               & Gestión de grupos y permisos. \textit{(solo staff)} \\
    Perfil         & Perfil               & Edición de los datos del usuario autenticado. \\
    \bottomrule
\end{tabularx}
\end{table}

Las secciones marcadas como \textit{solo staff} no aparecen en el menú para usuarios sin privilegios de staff.

\subsection{Tema claro / oscuro}

En la parte inferior de la barra lateral hay un botón de alternancia de tema. Haga clic en él para cambiar entre el \textbf{tema claro} y el \textbf{tema oscuro}. La preferencia se aplica de inmediato.

% ==============================================================
\section{Dashboard}
% ==============================================================

El Dashboard es la pantalla principal de la aplicación. Muestra el estado de los dispositivos registrados en el aula seleccionada y permite realizar acciones sobre ellos en tiempo real.

\subsection{Selección de aula}

En la parte superior del Dashboard encontrará un desplegable para seleccionar el \textbf{aula activa}. Al cambiar de aula, la lista de dispositivos se actualiza automáticamente.

\begin{infonote}
\textbf{Nota:} Los usuarios estándar solo pueden ver las aulas que tienen asignadas. Los superusuarios y el staff pueden acceder a cualquier aula.
\end{infonote}

\subsection{Control de internet del aula}

Junto al selector de aula aparece el botón \textbf{Bloquear internet del aula} / \textbf{Desbloquear internet del aula}. Este botón actúa sobre todos los dispositivos del aula simultáneamente:

\begin{itemize}
    \item \textbf{Bloquear (rojo):} Aplica las reglas de filtrado en el switch/router MikroTik para cortar el acceso a internet de todos los equipos del aula, excepto los dominios o IPs incluidos en la lista de Hosts Permitidos.
    \item \textbf{Desbloquear (verde):} Elimina las reglas de bloqueo y restaura el acceso total.
\end{itemize}

El estado del botón se actualiza automáticamente vía WebSocket: si otro usuario realiza cambios, el dashboard refleja el nuevo estado sin necesidad de recargar la página.

\begin{tipnote}
\textbf{Consejo:} El botón aparece en rojo si \textit{todos} los dispositivos del aula están bloqueados, y en verde si alguno tiene acceso libre. Puede usar el botón global para garantizar un estado uniforme en toda el aula.
\end{tipnote}

\subsection{Tarjetas de dispositivo}

Cada dispositivo del aula se representa con una \textbf{tarjeta} que muestra:

\begin{itemize}
    \item \textbf{Nombre del equipo} (hostname).
    \item \textbf{Dirección IP} asignada.
    \item \textbf{Sistema operativo} detectado.
    \item \textbf{Indicador de estado en línea} (punto verde = conectado, gris = desconectado).
    \item \textbf{Indicador de internet} (icono de globo = con acceso; icono tachado = bloqueado).
    \item \textbf{Switch/puerto} al que está conectado físicamente (si está configurado).
\end{itemize}

\subsubsection{Panel de acciones del dispositivo}

Al hacer clic sobre una tarjeta de dispositivo se abre un \textbf{panel inferior} con las siguientes acciones:

\begin{table}[H]
\centering
\begin{tabularx}{\textwidth}{lX}
    \toprule
    \textbf{Acción} & \textbf{Descripción} \\
    \midrule
    \textbf{Bloquear / Desbloquear internet} & Activa o desactiva el acceso a internet solo para ese equipo. \\
    \textbf{Apagar equipo} & Envía una orden de apagado al agente cliente instalado en el equipo. Solicita confirmación antes de ejecutarse. \\
    \textbf{Ver captura de pantalla} & Muestra la última captura de pantalla recibida del equipo (si el agente la ha enviado). Las capturas se conservan durante la sesión del navegador. \\
    \textbf{Eliminar dispositivo} & Elimina el registro del dispositivo del sistema. Solicita confirmación. Solo disponible para staff. \\
    \bottomrule
\end{tabularx}
\end{table}

\begin{warnnote}
\textbf{Atención:} La acción de apagado es inmediata una vez confirmada. Asegúrese de que los alumnos han guardado su trabajo antes de ejecutarla.
\end{warnnote}

\subsection{Paginación}

Si el aula tiene más de 10 dispositivos, la lista se divide en páginas. Use los controles de paginación en la parte inferior para navegar entre ellas.

% ==============================================================
\section{Gestión de aulas}
% ==============================================================

\textit{Esta sección requiere permiso de staff y el permiso \texttt{view\_classroom}.}

Acceda a \textbf{Classroom} desde el menú lateral para gestionar las aulas del centro.

\subsection{Listado de aulas}

Se muestra una lista paginada (10 elementos por página) con las aulas registradas. Cada tarjeta de aula indica:

\begin{itemize}
    \item Nombre del aula.
    \item Número de dispositivos asignados.
\end{itemize}

Use el campo de \textbf{búsqueda} en la parte superior para filtrar aulas por nombre.

\subsection{Crear una nueva aula}

\begin{enumerate}
    \item Haga clic en el botón \textbf{+ Nueva aula}.
    \item En el modal que se abre, introduzca el \textbf{nombre} del aula.
    \item Seleccione los \textbf{dispositivos} que desea asignar al aula (campo de selección múltiple).
    \item Pulse \textbf{Guardar}.
\end{enumerate}

\subsection{Editar un aula}

Haga clic en el icono de \textbf{edición} (lápiz) de la tarjeta correspondiente. Se abrirá el mismo modal con los datos actuales del aula. Modifique los campos necesarios y pulse \textbf{Guardar}.

\subsection{Eliminar un aula}

Haga clic en el icono de \textbf{papelera} de la tarjeta. Se solicitará confirmación. Una vez confirmado, el aula se elimina y los dispositivos asignados quedan sin aula.

\begin{warnnote}
\textbf{Atención:} Eliminar un aula no elimina los dispositivos que tenía asignados. Estos quedarán disponibles para ser reasignados.
\end{warnnote}

% ==============================================================
\section{Gestión de switches y routers}
% ==============================================================

\textit{Esta sección requiere permiso de staff y el permiso \texttt{view\_networkdevice}.}

Acceda a \textbf{Switches / Routers} desde el menú lateral para registrar y gestionar los equipos de red MikroTik que controlan el acceso a internet de las aulas.

\subsection{Listado de equipos de red}

Cada tarjeta muestra:

\begin{itemize}
    \item Nombre descriptivo del equipo.
    \item Dirección IP de gestión.
    \item Usuario de conexión a la API MikroTik.
\end{itemize}

\subsection{Añadir un equipo de red}

\begin{enumerate}
    \item Pulse \textbf{+ Nuevo equipo}.
    \item Rellene los campos del formulario:
    \begin{itemize}
        \item \textbf{Nombre:} Identificador descriptivo (p. ej., «Switch Aula 1»).
        \item \textbf{IP de gestión:} Dirección IP del equipo MikroTik.
        \item \textbf{Usuario API:} Nombre de usuario con acceso a la API RouterOS.
        \item \textbf{Contraseña API:} Contraseña del usuario API.
    \end{itemize}
    \item Pulse \textbf{Guardar}.
\end{enumerate}

\subsection{Editar o eliminar un equipo de red}

Use los iconos de \textbf{edición} o \textbf{papelera} en cada tarjeta para modificar los datos o eliminar el equipo del sistema.

\begin{warnnote}
\textbf{Atención:} Si elimina un equipo de red asociado a dispositivos activos, las acciones de bloqueo de internet sobre esos dispositivos dejarán de funcionar.
\end{warnnote}

% ==============================================================
\section{Hosts permitidos}
% ==============================================================

\textit{Esta sección requiere permiso de staff y el permiso \texttt{view\_allowedhost}.}

Los \textbf{hosts permitidos} son direcciones IP o dominios que permanecen accesibles para los alumnos \textit{incluso cuando el internet del aula está bloqueado}. Esto permite, por ejemplo, mantener el acceso a un servidor local de recursos educativos.

\subsection{Listado de hosts permitidos}

Cada tarjeta muestra la dirección IP o dominio, y si el host está restringido a un aula concreta o es \textbf{global} (accesible desde cualquier aula).

\subsection{Añadir un host permitido}

\begin{enumerate}
    \item Pulse \textbf{+ Nuevo host}.
    \item Complete el formulario:
    \begin{itemize}
        \item \textbf{Dirección:} IP o dominio a permitir (p. ej., \texttt{192.168.1.50} o \texttt{moodle.centro.es}).
        \item \textbf{Aula} (opcional): Si se selecciona un aula, el host solo estará permitido para los dispositivos de esa aula. Si se deja vacío, el host es \textbf{global}.
    \end{itemize}
    \item Pulse \textbf{Guardar}.
\end{enumerate}

\begin{tipnote}
\textbf{Consejo:} Use hosts globales para recursos comunes (servidor de ficheros, impresoras) y hosts de aula para recursos específicos de un grupo (servidor de prácticas de una asignatura concreta).
\end{tipnote}

\subsection{Editar o eliminar un host}

Use los iconos de \textbf{edición} o \textbf{papelera} en cada tarjeta.

% ==============================================================
\section{Gestión de usuarios}
% ==============================================================

\textit{Esta sección requiere permiso de staff y el permiso \texttt{view\_user}.}

\subsection{Listado de usuarios}

Se muestra una lista paginada con todos los usuarios del sistema. Use el buscador en la parte superior para filtrar por nombre de usuario.

Cada tarjeta de usuario muestra:
\begin{itemize}
    \item Avatar con la inicial del usuario (dorado para superusuarios).
    \item Nombre completo y nombre de usuario.
    \item Grupos y roles asignados.
    \item Aulas asignadas al usuario.
    \item Botones de \textbf{edición} y \textbf{eliminación}.
\end{itemize}

\subsection{Crear un usuario}

\begin{enumerate}
    \item Pulse \textbf{+ Nuevo usuario}.
    \item Rellene el formulario con los datos del usuario:
    \begin{itemize}
        \item \textbf{Nombre de usuario} (obligatorio, único en el sistema).
        \item \textbf{Nombre} y \textbf{apellidos}.
        \item \textbf{Correo electrónico}.
        \item \textbf{Contraseña} y confirmación.
        \item \textbf{Grupos:} Asigne uno o varios grupos de permisos.
        \item \textbf{Aulas:} Aulas a las que tendrá acceso el usuario.
        \item \textbf{Staff / Superusuario:} Marque si corresponde.
    \end{itemize}
    \item Pulse \textbf{Guardar}.
\end{enumerate}

\subsection{Editar un usuario}

Haga clic en el icono de \textbf{edición} de la tarjeta. Se abrirá el modal con los datos actuales. El campo de contraseña solo es necesario si desea cambiarla; si se deja vacío, se conserva la contraseña anterior.

\subsection{Eliminar un usuario}

Haga clic en el icono de \textbf{papelera}. Se solicitará confirmación. La eliminación es permanente.

% ==============================================================
\section{Gestión de grupos y permisos}
% ==============================================================

\textit{Esta sección requiere permiso de staff y el permiso \texttt{view\_group}.}

Los grupos permiten agrupar permisos y asignarlos a varios usuarios a la vez. Cada permiso controla el acceso a una operación concreta (ver, crear, modificar, eliminar) sobre una entidad del sistema.

\subsection{Listado de grupos}

Se muestra una lista paginada. Cada tarjeta de grupo indica su nombre y los permisos que tiene asignados como \textit{chips} (etiquetas).

\subsection{Crear un grupo}

\begin{enumerate}
    \item Pulse \textbf{+ Nuevo grupo}.
    \item Introduzca el \textbf{nombre} del grupo.
    \item Use el menú desplegable \textbf{Añadir permisos} para seleccionar los permisos que tendrá el grupo.
    \item Pulse \textbf{Guardar}.
\end{enumerate}

\subsection{Editar un grupo}

Haga clic en el icono de \textbf{edición} de la tarjeta. Puede:
\begin{itemize}
    \item Cambiar el nombre del grupo.
    \item Añadir nuevos permisos desde el menú desplegable.
    \item Eliminar permisos existentes haciendo clic en la \textbf{X} de cada chip de permiso.
\end{itemize}

\subsection{Eliminar un grupo}

Haga clic en el icono de \textbf{papelera} dentro del modal de edición. Se solicitará confirmación. Los usuarios que pertenecían a ese grupo conservan su cuenta, pero pierden los permisos del grupo eliminado.

\subsection{Referencia de permisos}

\begin{table}[H]
\centering
\begin{tabularx}{\textwidth}{lX}
    \toprule
    \textbf{Permiso} & \textbf{Acción que permite} \\
    \midrule
    \texttt{view\_device}        & Ver dispositivos en el Dashboard. \\
    \texttt{view\_classroom}     & Acceder a la sección de aulas. \\
    \texttt{add/change/delete\_classroom} & Crear, editar o eliminar aulas. \\
    \texttt{view\_networkdevice} & Ver equipos de red. \\
    \texttt{add/change/delete\_networkdevice} & Gestionar equipos de red. \\
    \texttt{view\_allowedhost}   & Ver hosts permitidos. \\
    \texttt{add/change/delete\_allowedhost} & Gestionar hosts permitidos. \\
    \texttt{view\_user}          & Ver la lista de usuarios. \\
    \texttt{add/change/delete\_user} & Gestionar usuarios. \\
    \texttt{view\_group}         & Ver grupos de permisos. \\
    \texttt{add/change/delete\_group} & Gestionar grupos. \\
    \bottomrule
\end{tabularx}
\end{table}

% ==============================================================
\section{Perfil de usuario}
% ==============================================================

Acceda a \textbf{Perfil} desde el menú lateral para consultar y modificar los datos de su propia cuenta.

\subsection{Información mostrada}

\begin{itemize}
    \item Avatar con la inicial del nombre.
    \item Nombre de usuario, nombre completo y correo electrónico.
    \item Roles activos: se muestran \textit{chips} indicando si el usuario es \textbf{Superusuario}, \textbf{Staff} o \textbf{Usuario activo}.
    \item Grupos a los que pertenece.
\end{itemize}

\subsection{Editar datos personales}

\begin{enumerate}
    \item Modifique los campos de \textbf{nombre}, \textbf{apellidos} o \textbf{correo electrónico}.
    \item Pulse el botón \textbf{Guardar cambios} (icono de disco).
\end{enumerate}

\subsection{Cambiar la contraseña}

\begin{enumerate}
    \item Introduzca la \textbf{nueva contraseña} en el campo correspondiente.
    \item Repita la contraseña en el campo de \textbf{confirmación}.
    \item Pulse \textbf{Guardar cambios}.
\end{enumerate}

\begin{warnnote}
\textbf{Atención:} Si introduce una contraseña nueva y los dos campos no coinciden, el sistema mostrará un error y los cambios no se guardarán.
\end{warnnote}

\begin{infonote}
\textbf{Nota:} Si no desea cambiar la contraseña, deje ambos campos en blanco. Se conservará la contraseña actual.
\end{infonote}

% ==============================================================
\section{Preguntas frecuentes}
% ==============================================================

\begin{description}[style=nextline]

\item[\textbf{El Dashboard no muestra ningún dispositivo.}]
Compruebe que hay un aula seleccionada en el desplegable. Si el problema persiste, verifique que los agentes cliente están en funcionamiento en los equipos del aula.

\item[\textbf{Un dispositivo aparece como desconectado aunque el equipo está encendido.}]
El agente cliente puede haberse detenido o el equipo puede haber perdido conectividad de red. Compruebe que el servicio \texttt{netmanagement-client} está activo en el equipo afectado.

\item[\textbf{El botón de bloquear internet no funciona.}]
Verifique que hay al menos un equipo de red (switch/router MikroTik) configurado en la sección \textit{Switches / Routers} y que la IP de gestión es accesible desde el servidor.

\item[\textbf{No veo la sección «Aulas» ni «Usuarios» en el menú.}]
Estas secciones requieren privilegios de \textbf{staff}. Contacte con el administrador del sistema para que le asigne el rol o los permisos necesarios.

\item[\textbf{He olvidado mi contraseña.}]
Contacte con un administrador del sistema para que le restablezca la contraseña desde la sección de \textit{Gestión de usuarios}.

\item[\textbf{La captura de pantalla de un dispositivo es antigua.}]
Las capturas se actualizan periódicamente según la configuración del agente cliente. Si la imagen es muy antigua, es posible que el agente esté detenido.

\end{description}

\end{document}
